% LaTeX Template for AI Problem Analysis Output
% AMC12A 2018 Problem 15 - Complete Strategic Analysis
% IMPORTANT: Use LaTeX commands for all mathematical symbols, NOT unicode characters

\documentclass[]{article}
\usepackage[utf8]{inputenc}
\usepackage{amsmath, amssymb, amsthm}
\usepackage{geometry}
\usepackage{xcolor}
\usepackage[most]{tcolorbox}
\usepackage{enumitem}
\usepackage{fancyhdr}
\usepackage{asymptote}

% Page setup
\geometry{margin=1in}
\pagestyle{fancy}
\fancyhf{}
\rhead{AMC12/COMC Problem Analysis}
\lhead{\today}

% Color definitions
\definecolor{problemconfig}{RGB}{230, 242, 255}
\definecolor{strategic}{RGB}{255, 230, 242}
\definecolor{tactical}{RGB}{242, 255, 230}
\definecolor{techniques}{RGB}{255, 242, 230}
\definecolor{verification}{RGB}{255, 230, 230}
\definecolor{solution}{RGB}{230, 255, 255}
\definecolor{competition}{RGB}{242, 242, 255}
\definecolor{gold}{RGB}{218, 165, 32}

% Box styles
\newtcolorbox{problemconfigbox}{
    colback=problemconfig,
    colframe=blue,
    title=Problem Configuration Analysis,
    fonttitle=\bfseries,
    arc=3pt,
    boxrule=1pt,
    breakable
}

\newtcolorbox{strategicbox}{
    colback=strategic,
    colframe=purple,
    title=Strategic Framework Application,
    fonttitle=\bfseries,
    arc=3pt,
    boxrule=1pt,
    breakable
}

\newtcolorbox{tacticalbox}{
    colback=tactical,
    colframe=green,
    title=Tactical Methodology Selection,
    fonttitle=\bfseries,
    arc=3pt,
    boxrule=1pt,
    breakable
}

\newtcolorbox{techniquesbox}{
    colback=techniques,
    colframe=orange,
    title=Conversion Techniques Application,
    fonttitle=\bfseries,
    arc=3pt,
    boxrule=1pt,
    breakable
}

\newtcolorbox{verificationbox}{
    colback=verification,
    colframe=red,
    title=Verification Strategy Design,
    fonttitle=\bfseries,
    arc=3pt,
    boxrule=1pt,
    breakable
}

\newtcolorbox{solutionbox}{
    colback=solution,
    colframe=cyan,
    title=Solution Roadmap,
    fonttitle=\bfseries,
    arc=3pt,
    boxrule=1pt,
    breakable
}

\newtcolorbox{competitionbox}{
    colback=competition,
    colframe=purple,
    title=Competition Strategy Integration,
    fonttitle=\bfseries,
    arc=3pt,
    boxrule=1pt,
    breakable
}

\newtcolorbox{keyinsightbox}{
    colback=yellow!20,
    colframe=gold,
    title=Key Insight,
    fonttitle=\bfseries,
    arc=3pt,
    boxrule=2pt,
    leftrule=5pt,
    breakable
}

\newtcolorbox{warningbox}{
    colback=red!10,
    colframe=red!50,
    title=Warning: Common Pitfall,
    fonttitle=\bfseries,
    arc=3pt,
    boxrule=2pt,
    leftrule=5pt,
    breakable
}

\newtcolorbox{tipbox}{
    colback=green!10,
    colframe=green!50,
    title=Pro Tip,
    fonttitle=\bfseries,
    arc=3pt,
    boxrule=2pt,
    leftrule=5pt,
    breakable
}

\begin{document}

\title{AMC12A 2018 Problem 15: Strategic Problem Analysis}
\author{Competition Math Framework}
\date{\today}
\maketitle

\section*{Problem Statement}

A scanning code consists of a $7 \times 7$ grid of squares, with some of its squares colored black and the rest colored white. There must be at least one square of each color in this grid of $49$ squares. A scanning code is called \textit{symmetric} if its look does not change when the entire square is rotated by a multiple of $90^{\circ}$ counterclockwise around its center, nor when it is reflected across a line joining opposite corners or a line joining midpoints of opposite sides. What is the total number of possible symmetric scanning codes?

$$\textbf{(A) } 510 \qquad \textbf{(B) } 1022 \qquad \textbf{(C) } 8190 \qquad \textbf{(D) } 8192 \qquad \textbf{(E) } 65,534$$

\begin{problemconfigbox}
\subsection*{A. Problem Classification}
\begin{itemize}[leftmargin=*]
    \item \textbf{Problem Type}: To Find - count the number of distinct symmetric colorings
    \item \textbf{Difficulty Band}: Mid-band (Problem 15 of 25) - requires symmetry analysis and counting technique, but not advanced olympiad methods
    \item \textbf{Competition Context}: AMC12 (75 min, 25 problems) - approximately 3 minutes allocated; also appears as AMC 10A Problem 19
    \item \textbf{Mathematical Domain}: Combinatorics (counting with symmetry constraints), with connections to group theory (Burnside's Lemma) and geometry (dihedral symmetry)
\end{itemize}

\subsection*{B. Problem Structure Identification}
\begin{itemize}[leftmargin=*]
    \item \textbf{Given Information}: 
    \begin{itemize}
        \item $7 \times 7$ grid of squares (49 squares total)
        \item Each square colored black or white
        \item Symmetry group: dihedral group $D_4$ (4 rotations + 4 reflections)
        \item Rotational symmetry: invariant under $90^{\circ}$, $180^{\circ}$, $270^{\circ}$ rotations
        \item Reflectional symmetry: invariant under 2 diagonal reflections and 2 axis reflections
    \end{itemize}
    \item \textbf{Constraints}: 
    \begin{itemize}
        \item Must have at least one black square AND at least one white square
        \item Coloring must be invariant under all 8 symmetries of the square
    \end{itemize}
    \item \textbf{Objective}: Count the total number of valid symmetric colorings
    \item \textbf{Hidden Assumptions}: 
    \begin{itemize}
        \item The symmetry constraints drastically reduce the number of independent choices
        \item Many squares' colors are determined by symmetry once a few are chosen
        \item The center square is fixed by all symmetries
    \end{itemize}
    \item \textbf{Answer Choice Leverage}: 
    \begin{itemize}
        \item Options follow the pattern: $510 = 2^9 - 2$, $1022 = 2^{10} - 2$, $8190 = 2^{13} - 2$, $8192 = 2^{13}$, $65534 = 2^{16} - 2$
        \item The pattern $2^n - 2$ suggests we're counting binary choices minus two monochromatic cases
        \item This tells us to find the number of independent color choices
    \end{itemize}
\end{itemize}

\subsection*{C. Strategic Orientation}
\begin{itemize}[leftmargin=*]
    \item \textbf{Hypothesis}: The 8-fold symmetry divides the grid into equivalence classes
    \item \textbf{Conclusion}: Find the number of "fundamental domains" (independent positions), compute $2^n$ for $n$ independent choices, then subtract 2
    \item \textbf{Notation}: Label squares by their symmetry class (positions that map to each other under the symmetry group)
    \item \textbf{Intuitive Approach}: Visualize folding the grid along symmetry axes
    \item \textbf{Cross-Domain Potential}: Group Theory (Burnside's Lemma), Fundamental Domain Analysis
\end{itemize}
\end{problemconfigbox}

\begin{strategicbox}
\subsection*{A. Psychological Strategies}
\begin{itemize}[leftmargin=*]
    \item \textbf{Mental Toughness}: Don't be intimidated by the $7 \times 7$ grid - symmetry simplifies dramatically
    \item \textbf{Creativity Cultivation}: Try the simpler $3 \times 3$ case first to see the pattern
    \item \textbf{Iterative Refinement}: If first labeling is confusing, try a different scheme
\end{itemize}

\subsection*{B. Investigation Strategies}
\begin{itemize}[leftmargin=*]
    \item \textbf{Get Your Hands Dirty}: Draw the $7 \times 7$ grid and mark symmetry lines
    \item \textbf{Penultimate Step}: Answer will be $2^n - 2$ where $n$ is the number of independent regions
    \item \textbf{Wishful Thinking}: "I wish I only had to color 10 squares and the rest would be determined"
    \item \textbf{Make It Easier}: Try $3 \times 3$ grid: center (1), corner class (1), edge class (1) $\Rightarrow$ $2^3 - 2 = 6$
    \item \textbf{Visualization}: Systematically label squares from center outward
\end{itemize}

\subsection*{C. Late-Band Strategic Playbook}
\begin{itemize}[leftmargin=*]
    \item \textbf{Answer-Choice Leverage}: Pattern $2^{10} - 2 = 1022$ suggests $n = 10$ independent regions
    \item \textbf{Penultimate Step}: Work backwards - if answer is $2^{10} - 2$, there must be exactly 10 independent squares
\end{itemize}

\subsection*{D. Cross-Domain Translation}
\begin{itemize}[leftmargin=*]
    \item \textbf{Draw a Picture}: Essential - draw grid and mark symmetries
    \item \textbf{Recast the Problem}: "Count colorings" becomes "count independent binary choices"
    \item \textbf{Change Your Point of View}: View grid as concentric layers from center
\end{itemize}
\end{strategicbox}

\begin{tacticalbox}
\subsection*{A. Core Mathematical Tactics}
\begin{itemize}[leftmargin=*]
    \item \textbf{Symmetry (PRIMARY TACTIC)}: 
    \begin{itemize}
        \item Dihedral group $D_4$ has 8 elements: identity + 3 rotations + 4 reflections
        \item These partition the 49 squares into equivalence classes
        \item All squares in same class must have same color
        \item Count equivalence classes = count independent choices
    \end{itemize}
    \item \textbf{Fundamental Domain}: Identify minimal region that generates all squares
    \item \textbf{Binary Choice Counting}: $n$ independent regions $\Rightarrow$ $2^n$ total combinations
\end{itemize}

\subsection*{B. Domain-Specific Tactics}
\begin{itemize}[leftmargin=*]
    \item \textbf{Combinatorics}: Multiplication principle ($2^n$) and subtraction principle ($-2$)
    \item \textbf{Symmetry Analysis}:
    \begin{itemize}
        \item Center square: fixed by all symmetries (1 independent position)
        \item Squares on axes: have 2-fold or 4-fold constraint
        \item General squares: form 8-element orbits
    \end{itemize}
    \item \textbf{Systematic Labeling}: Label from center outward using letters A, B, C, ...
\end{itemize}

\subsection*{C. Crossover Tactics}
\begin{itemize}[leftmargin=*]
    \item \textbf{Group Theory}: Apply Burnside's Lemma for rigorous counting
    \item \textbf{Geometry}: Visualize 8 symmetries of a square
    \item \textbf{Algebraic Coding}: Map colorings to binary strings of length $n$
\end{itemize}
\end{tacticalbox}

\begin{techniquesbox}
\subsection*{A. High-Probability Techniques}

\begin{itemize}[leftmargin=*]
    \item \textbf{Fundamental Domain Analysis (100\% applicable)}
    \begin{itemize}
        \item \textbf{Recognition}: Symmetries partition the domain
        \item \textbf{Why It Works}: Once colors chosen in fundamental domain, symmetry determines rest
        \item \textbf{Application}:
        \begin{enumerate}
            \item Draw $7 \times 7$ grid
            \item Mark 4 symmetry lines (2 diagonals, 2 axes)
            \item Label squares within one sector
            \item Count distinct labels = independent choices
        \end{enumerate}
    \end{itemize}
    
    \item \textbf{Burnside's Lemma / P\'olya Counting (TM2 - 95\% applicable)}
    \begin{itemize}
        \item \textbf{Recognition}: Count colorings with rotational/reflectional symmetry
        \item \textbf{Why It Works}: Dihedral group $D_4$ acts on 49 squares
        \item \textbf{Application}: For each symmetry, count fixed colorings and average
        \item \textbf{Note}: Theoretically rigorous but slower for competition
    \end{itemize}
\end{itemize}

\subsection*{B. Technique Integration}

\textbf{Primary Technique}: Fundamental Domain with Systematic Labeling

\textbf{Why This Works Best}:
\begin{enumerate}
    \item Speed: 2-3 minutes with practice
    \item Accuracy: Visual and systematic
    \item Competition-Friendly: No advanced group theory needed
    \item Verification-Ready: Can check multiple ways
\end{enumerate}

\textbf{Integration Strategy}:
\begin{enumerate}
    \item Draw $7 \times 7$ grid (15 sec)
    \item Mark 4 symmetry lines (10 sec)
    \item Label from center outward (60 sec)
    \item Verify symmetry (30 sec)
    \item Count letters: 10 (15 sec)
    \item Compute: $2^{10} - 2 = 1022$ (15 sec)
    \item Match answer (5 sec)
\end{enumerate}
\end{techniquesbox}

\begin{verificationbox}
\subsection*{A. Verification Level Selection}

\textbf{Selected Level}: \textbf{Level 4 Enhanced} (standard with alternative method)

\textbf{Rationale}: Mid-band problem (P15), high confidence expected, 30-45 seconds available

\subsection*{B. Verification Methods}

\textbf{Primary (20 seconds)}:
\begin{enumerate}[nosep]
    \item Sanity check: 1022 reasonable? (49 squares $\rightarrow$ huge without symmetry, 1000 with 8-fold symmetry seems right)
    \item Pattern check: $1022 = 2^{10} - 2$ confirms 10 independent regions
    \item Small case: $3 \times 3$ gives $2^3 - 2 = 6$ (mental check)
\end{enumerate}

\textbf{Alternative Method (25 seconds)}:
\begin{enumerate}[nosep]
    \item Folding visualization: fold along 4 lines, count non-overlapping squares
    \item Layer count: center (1) + layer 1 (3 types) + layer 2 (3 types) + outer (3 types) = 10
\end{enumerate}

\textbf{Boundary Cases (10 seconds)}:
\begin{itemize}[nosep]
    \item All black: 1 config (excluded)
    \item All white: 1 config (excluded)
    \item One black at center: 1 config (included)
\end{itemize}

\subsection*{C. Confidence Calibration}

\begin{itemize}[leftmargin=*]
    \item Initial: 90\% (careful labeling)
    \item After primary: 95\% (pattern matches)
    \item After alternative: 98\% (two methods agree)
    \item \textbf{Final: 98\%} - Very high, proceed
\end{itemize}
\end{verificationbox}

\begin{solutionbox}
\subsection*{A. Step-by-Step Solution}

\textbf{Step 1: Understand Symmetries (30 sec)}

The code must be invariant under:
\begin{itemize}[nosep]
    \item 4 rotations: $0^{\circ}, 90^{\circ}, 180^{\circ}, 270^{\circ}$
    \item 4 reflections: 2 diagonals + 2 midpoint axes
\end{itemize}

\textbf{Step 2: Label the Grid (90 sec)}

$$\begin{array}{|c|c|c|c|c|c|c|}
\hline
K & J & H & G & H & J & K \\
\hline
J & F & E & D & E & F & J \\
\hline
H & E & C & B & C & E & H \\
\hline
G & D & B & A & B & D & G \\
\hline
H & E & C & B & C & E & H \\
\hline
J & F & E & D & E & F & J \\
\hline
K & J & H & G & H & J & K \\
\hline
\end{array}$$

\textbf{Note}: Letter I skipped to avoid confusion with number 1

\textbf{Step 3: Count Labels (15 sec)}

Letters used: A, B, C, D, E, F, G, H, J, K = \textbf{10 distinct labels}

\textbf{Step 4: Apply Formula (15 sec)}

$$\text{Total} = 2^{10} = 1024$$
$$\text{Valid} = 2^{10} - 2 = 1022$$

\textbf{Step 5: Match Answer (5 sec)}

$$\boxed{\textbf{(B) } 1022}$$

\subsection*{B. Alternative: Folding Method}

\begin{asy}
size(100pt);
draw((1,0)--(8,0),linewidth(0.5));
draw((1,2)--(6,2),linewidth(0.5));
draw((1,4)--(4,4),linewidth(0.5));
draw((1,6)--(2,6),linewidth(0.5));
draw((2,6)--(2,0),linewidth(0.5));
draw((4,4)--(4,0),linewidth(0.5));
draw((6,2)--(6,0),linewidth(0.5));
draw((1,0)--(1,7),dashed+linewidth(0.5));
draw((1,7)--(8,0),dashed+linewidth(0.5));
\end{asy}

Fold grid along all 4 symmetry lines. Non-overlapping squares = 10 independent regions.

Therefore: $2^{10} - 2 = 1022$

\subsection*{C. Alternative: Corner Squares + Stripes}

Divide the $7 \times 7$ grid into:
\begin{itemize}[nosep]
    \item Four $3 \times 3$ corner squares
    \item Horizontal and vertical stripes through middle
\end{itemize}

For one $3 \times 3$ corner:
$$\begin{array}{|c|c|c|}
\hline
A & B & C \\
\hline
B & D & E \\
\hline
C & E & F \\
\hline
\end{array}$$

6 letters $\Rightarrow$ $2^6 = 64$ colorings

For stripes: 4 distinct positions $\Rightarrow$ $2^4 = 16$ colorings

Total: $64 \times 16 = 1024$

Subtract monochromatic: $1024 - 2 = 1022$

\subsection*{D. Time Management}

\begin{itemize}[leftmargin=*]
    \item Reading and setup: 30 sec
    \item Labeling: 90 sec
    \item Counting and calculation: 30 sec
    \item Verification: 30 sec
    \item \textbf{Total: 3 minutes}
\end{itemize}

\subsection*{E. Pivot Indicators}

\textbf{Switch if}:
\begin{itemize}[nosep]
    \item Labeling confusing after 1 min $\Rightarrow$ try folding
    \item Can't see 10 regions $\Rightarrow$ try layer-by-layer
    \item Arithmetic wrong $\Rightarrow$ check $2^n - 2$ pattern
\end{itemize}

\textbf{Abandon if}:
\begin{itemize}[nosep]
    \item No progress after 4 minutes
    \item Getting 13 or 16 regions (wrong counting)
\end{itemize}
\end{solutionbox}

\begin{competitionbox}
\subsection*{A. Time Management}

\textbf{Allocated}: 3 minutes (mid-band)

\textbf{Time-Saving}:
\begin{itemize}[nosep]
    \item Rough sketch sufficient (don't draw neatly)
    \item Use folding if comfortable (saves 30-60 sec)
    \item Recognize $2^n - 2$ pattern immediately
\end{itemize}

\subsection*{B. Problem Prioritization}

\textbf{Attempt if}:
\begin{itemize}[nosep]
    \item Comfortable with symmetry problems
    \item Aiming for 15+ correct
    \item Secured easier problems (P1-14)
\end{itemize}

\textbf{Skip if}:
\begin{itemize}[nosep]
    \item Struggling with P10-14
    \item Running low on time
\end{itemize}

\subsection*{C. Risk Management}

\textbf{Confidence Thresholds}:
\begin{itemize}[leftmargin=*]
    \item \textbf{High (90\%+)}: Submit confidently
    \item \textbf{Medium (70-90\%)}: Verify with folding, then submit
    \item \textbf{Low (50-70\%)}: Check if count matches answer pattern
    \item \textbf{Very Low (<50\%)}: Educated guess
\end{itemize}

\textbf{Error Recovery}:
\begin{itemize}[nosep]
    \item 13 regions: didn't account for all symmetries
    \item 16 regions: counted dependent squares as independent
    \item 9 regions: merged distinct classes incorrectly
\end{itemize}

\textbf{Answer Analysis}:
\begin{itemize}[nosep]
    \item (A) $510 = 2^9 - 2$: 9 regions
    \item (B) $1022 = 2^{10} - 2$: 10 regions $\checkmark$
    \item (C) $8190 = 2^{13} - 2$: 13 regions
    \item (D) $8192 = 2^{13}$: forgot to subtract 2
    \item (E) $65534 = 2^{16} - 2$: 16 regions (too many)
\end{itemize}

\subsection*{D. Optimization}

\textbf{Expected Value}: $0.75 \times 6 = 4.5$ pts vs 1.5 for blank

\textbf{Strategic Guessing}: If must guess, eliminate (D) and (E), choose (B) as middle option

\subsection*{E. Pattern Recognition}

\textbf{Similar Problems}: 1996 AIME \#7

\textbf{Pattern}: Colorings with symmetry $\Rightarrow$ fundamental domain $\Rightarrow$ $2^n - k$

\textbf{Abandonment}:
\begin{itemize}[nosep]
    \item Soft: No progress after 5 min (may return)
    \item Hard: Confusion about symmetries after 2 min
\end{itemize}
\end{competitionbox}

\begin{keyinsightbox}
The 8 symmetries of dihedral group $D_4$ partition the $7 \times 7$ grid into exactly \textbf{10 equivalence classes}. By systematically labeling from the center outward (A, B, C, D, E, F, G, H, J, K), we identify these 10 independent color choices. Since each can be black or white: $2^{10} = 1024$ total, minus 2 monochromatic cases = $\boxed{1022}$ valid codes.

The answer pattern $2^n - 2$ immediately signals to find $n$ independent binary choices!
\end{keyinsightbox}

\begin{warningbox}
\textbf{Most Common Errors}:

\begin{enumerate}
    \item \textbf{Forgetting subtraction}: Getting $1024$ and selecting (D) $8192 = 2^{13}$ by misreading
    \item \textbf{Miscounting classes}: Getting 13 (not recognizing all symmetries) or 9 (merging incorrectly)
    \item \textbf{Symmetry confusion}: Misunderstanding diagonal reflections
    \item \textbf{Arithmetic errors}: Miscalculating $2^{10}$ or forgetting $-2$
\end{enumerate}

\textbf{Prevention}: Always verify labeling satisfies at least one $90^{\circ}$ rotation and one reflection!
\end{warningbox}

\begin{tipbox}
\textbf{Competition Strategy}: When you see "symmetric under rotations and reflections":
\begin{enumerate}[nosep]
    \item Draw symmetry lines
    \item Look for $2^n - k$ pattern in choices
    \item Count independent regions = $n$
    \item Verify $n$ gives an answer choice
\end{enumerate}

\textbf{Quick Check}: For $7 \times 7$ with 8-fold symmetry, 10 independent regions is reasonable: center (1) + 3 layers $\times$ 3 types each $\approx$ 10

\textbf{Time-Saver}: Folding visualization solves in under 2 minutes with practice!

\textbf{Related}: 1996 AIME \#7, 2015 AMC 10A \#20 (similar symmetry-counting)
\end{tipbox}

\section*{Final Answer}
$$\boxed{\textbf{(B) } 1022}$$

\section*{Confidence Assessment}
\begin{itemize}[leftmargin=*]
    \item \textbf{Confidence}: 98\%
    \item \textbf{Verification}: Level 4 Enhanced (multiple methods)
    \item \textbf{Flag for Review}: No
    \item \textbf{Time Spent}: 3 minutes (on target)
\end{itemize}

\section*{Solution Summary}

\textbf{Strategy}: Fundamental domain identification via systematic symmetry-based labeling

\textbf{Key Steps}:
\begin{enumerate}
    \item Recognize dihedral group $D_4$ (8 symmetries)
    \item Draw grid and mark 4 symmetry lines
    \item Label squares ensuring symmetry preservation
    \item Identify 10 distinct equivalence classes (A-K, skipping I)
    \item Apply: $2^{10} - 2 = 1022$
    \item Verify and match answer (B)
\end{enumerate}

\textbf{Why This Works}: Symmetry forces squares in same equivalence class to have same color. We have 10 independent binary choices, giving $2^{10} = 1024$ total colorings. Excluding all-black and all-white gives $\boxed{1022}$ valid symmetric scanning codes.

\end{document}